
Generative policies based on diffusion models
\cite{song2019generative,song2020score,ho2020denoising}
and flow matching models
\cite{lipman2022flow,liu2022flow}
have emerged as a powerful paradigm for policy learning.
Diffusion policies generate actions through iterative denoising, while flow policies generate actions by integrating a learned velocity field, often providing a deterministic and more efficient alternative.
Due to their flexibility and accuracy in modeling high-dimensional action distributions, generative policies have achieved strong empirical gains over deterministic or unimodal policies in challenging continuous-control settings
\cite{ke20243d,chi2025diffusion,huang2024diffuseloco,black2024pi,intelligence2025pi}.

A natural way to further improve generative policies is through reinforcement learning fine-tuning.
However, a central challenge is that the exact likelihoods of generative policies are generally intractable, since each action is produced only after simulating an entire denoising or flow trajectory.
This prevents the direct use of many stable and widely adopted policy optimization methods, such as PPO
\cite{schulman2017proximal}
and GRPO
\cite{shao2024deepseekmath}.
Several recent works have attempted to bypass this likelihood challenge
\cite{psenka2023learning,ding2024diffusion,ma2025efficient,ren2024diffusion,zhang2025reinflow,mcallister2025flow}.
However, these methods often suffer from biased estimation, limited data efficiency, substantial computational overhead, or restriction to a particular generative policy class or policy update rule.
Consequently, the full potential of RL for diffusion and flow policies remains underexplored.

To address this problem, we propose \textbf{Doob Policy Optimization (DoobPO)}, a likelihood-free framework that separates policy improvement from generative policy learning.
The first stage represents the desired policy improvement through a target policy ratio.
The second stage distills the resulting ratio-tilted target into the generative policy.
This decomposition avoids evaluating the likelihood of the generative policy itself.
The policy-improvement stage is agnostic to the underlying generative model and accommodates a broad family of updates, including regularized Policy Mirror Descent~\cite{zhan2023policy}, Soft Actor-Critic~\cite{haarnoja2018soft}, PPO~\cite{schulman2017proximal}, and GRPO~\cite{shao2024deepseekmath}.
The generative stage uses Doob's $h$-transform to characterize the ratio-tilted target process and learns this process through Doob's Drift Matching.
The same construction applies to both diffusion and flow policies.

\textbf{Contributions.}
Our contributions are three-fold:
\begin{itemize}
    \item \textbf{Method.}
    We propose \textbf{Doob Policy Optimization (DoobPO)}, a unified framework for RL fine-tuning of generative policies without evaluating exact action likelihoods
    (\cref{sec:method}).
    DoobPO follows a two-stage design: it first obtains the target policy ratio and then transfers the resulting policy improvement into the generative policy through a tractable conditional drift matching objective.
    The framework supports both diffusion and flow policies
    (\cref{sec:diff_dr_matching,sec:flow_dr_matching}).
    We summarize
    the procedures in
    \cref{alg:rwdm_ppo,alg:rwdm_ppo_fm}.
    Instantiating DoobPO with different policy updates yields
    \textbf{DoobPMD}, \textbf{DoobPPO}, and \textbf{DoobGRPO}.
    We further show that several existing methods arise as special cases under particular choices of the policy update and generative policy class
    (\cref{sec:unified_framework}).
    We summarize these instantiations in \Cref{tab:unified_framework}.

    \item \textbf{Theory.}
    We characterize the ratio-tilted target policy through Doob's $h$-transform
    (\cref{lem:algo_correctness}).
    This characterization leads to Doob's Drift Matching as a principled way to realize the desired policy update.
    For tractability, we introduce a conditional drift matching objective and prove that its gradient is proportional to that of the ideal Doob's Drift Matching objective for both diffusion and flow policies
    (\cref{thm:grad_equi_diff,thm:grad_equi_fm}).
    From the Stochastic Optimal Control perspective
    (\cref{sec:soc}),
    we further show that the Doob-transformed diffusion solves the corresponding endpoint-constrained control problem
    (\cref{prop:cdm_solves_soc}).
    Hence, DoobPO learns the minimum-energy modification of the pretrained diffusion process that realizes the terminal tilt
    (\cref{rem:ratio_h_interpretation}).

    \item \textbf{Experiments.}
    We evaluate DoobPO on MuJoCo continuous-control benchmarks with both diffusion and flow policies
    (\cref{sec:exp}).
    For diffusion policies, DoobPO achieves the strongest overall performance against competitive diffusion-policy RL baselines, with leading results on most tasks
    (\cref{sec:diffusion_exp}).
    For flow policies, DoobGRPO achieves the best performance on 9 out of 10 tasks under the same interaction budget and delivers large gains on challenging locomotion environments
    (\cref{sec:flow_exp}).
    These results establish DoobPO as a strong and practical framework for RL fine-tuning across both diffusion and flow policies.
\end{itemize}

\textbf{Notation.}
We summarize notation and time conventions in \cref{sec:notation}.

\textbf{Related Work.}
We defer an extended discussion of related work to \cref{sec:related_works} due to page limits.